% Drawing 85/02-1 -- folder 85, batch 02, archive page 33 (PDF page 34).
% Transcription: à gauche, à l'encre : deux arcs épais, l'un en haut (un petit « $0$ » au-dessus), l'autre plus bas ; au-dessous une petite fourche à trois branches, « $1$ » au bout de la branche du haut, « $0$ » sous le nœud, « $J$ » et « $J'$ » au bout des deux branches du bas
% Kind: two heavy arcs and a small three-branched fork with labels. Method 2 (schematic redraw from the facsimile). Status: partial.
% Notes: Ink. The upper heavy arc with a small ring above it and a tick and a dash at its right end; the lower heavy arc with a pencil ring at its upper end and a tick below; the fork: three branches from one node, 1 at the end of the upper one, 0 under the node, J' at the end of the right one; the hook at the end of the left branch is read J (the transcription's note has only 0, 1, J'). Left out: the long pencil lines that run down the left of the page and across block c); they cross the text and are taken as deletion strokes and a frame, not as part of a figure (the note leaves this open). Still differs: whether the two heavy arcs and the fork make one figure is not known; they are drawn at their places. Redrawn 2026-11-12 (issue 128): to check again.
% Coordinates: 1 unit = 1 mm = 1/100 of the crop width (bbox in page fractions: [0.09, 0.4, 0.3, 0.66]).
% Redrawn by Claude (claude-fable-5-2), issue #128, 2026-11-12; not checked by a human; not compiled (no TeX in the session).
% Rights: reproduces part of the work -- behind RENDER_TRANSCRIPTS, never published by itself.
\documentclass[tikz,border=3pt]{standalone}
\usepackage{amsmath,amssymb}
\usetikzlibrary{patterns,arrows.meta}
\tikzset{pencil/.style={black!45}, struck/.style={black!35}, mpt/.style={circle,draw,inner sep=1.3pt}, pt/.style={circle,fill,inner sep=0.9pt}, opt/.style={circle,draw,fill=white,inner sep=1.1pt}}
\begin{document}
\begin{tikzpicture}[x=1mm,y=1mm,line cap=round,line join=round,font=\scriptsize]
\draw[very thick] plot[smooth, tension=0.55] coordinates {(14,-38) (17,-33) (25,-25.5) (35,-20) (43,-19) (52,-21) (57,-24)};
\draw (58.5,-26) -- (59,-29.5);
\draw (33.5,-17.5) circle[radius=0.9];
\draw (54,-28) -- (56,-28);
\draw[very thick] plot[smooth, tension=0.55] coordinates {(25.5,-110.5) (26.5,-103) (31,-95) (36,-90.5) (40.5,-88)};
\draw[pencil] (45.5,-85.5) circle[radius=0.9];
\draw (22.5,-113) -- (22.5,-117);
\draw (30,-130.5) -- (37,-119);
\draw plot[smooth, tension=0.55] coordinates {(30,-130.5) (22,-139) (15.5,-150.5)};
\draw plot[smooth, tension=0.55] coordinates {(30,-130.5) (38,-134) (45.5,-140)};
\node[inner sep=1pt] at (41.3,-118) {\small $1$};
\node[inner sep=1pt] at (30,-135.7) {\small $0$};
\node[inner sep=1pt] at (48.5,-144) {\small $J'$};
\node[inner sep=1pt] at (14.5,-155) {\small $J$};
\end{tikzpicture}
\end{document}
